\paragraph{A strengthening with fixed edge constraints and linear Boolean support.}
Let $G=(V,E_G)$ belong to the bounded-degree expander family above,
and put $n=|V|$.  Increase the fixed degree bound to $d\ge3$
if necessary.  Retain only the two unary weights at each vertex:
\[
 J_G(z)=\sum_{I\subseteq V\,\text{independent}}
            \prod_{v\in I}z_{v,1}\prod_{v\notin I}z_{v,0}
                  \in\mathbb Z[z_{v,b}:v\in V, b\in\{0,1\}].
\]
This polynomial is homogeneous of degree $n$, and every monomial
contains exactly one variable from each vertex block
$\{z_{v,0},z_{v,1}\}$.  There is a constant $\gamma>0$, depending
only on the fixed expansion and degree bounds, such that every
binary positive arithmetic circuit over $\mathbb R_{\ge0}$ for
$J_G$ has at least $2^{\gamma n-2}$ internal gates for all
sufficiently large $n$.  Nevertheless NN-C034 supplies an integral
constant-free circuit of size and deterministic bit construction time
\[
 \operatorname{poly}_d(n)
 2^{O_d(n(\log\log n)^2/\sqrt{\log n})}=2^{o(n)},
\]
with polynomial formal degree.  The parameter is the original
number of vertices.  The exponential positive lower bound therefore
persists after every binary edge factor of the preceding universal
polynomial has been fixed to the independent-set constraint.
This assertion requires its own lower-bound argument; it does not
follow merely by specializing a polynomial with a lower bound.

For the same polynomial, Boolean evaluation of the coefficients and
monomials gives the allowed-weight support function
\[
 \operatorname{Supp}^{\rm ind}_G(a)
   =\bigvee_{I\,\text{independent}}
           \left(\bigwedge_{v\in I}a_{v,1}\wedge
                 \bigwedge_{v\notin I}a_{v,0}\right).
\]
It has the explicit monotone formula
\begin{equation}\label{eq:S019-independent-support}
 \operatorname{Supp}^{\rm ind}_G(a)=
       \bigwedge_{v\in V}(a_{v,0}\vee a_{v,1})
       \ \wedge\!
       \bigwedge_{uv\in E_G}(a_{u,0}\vee a_{v,0}).
\end{equation}
In particular it uses at most $2(n+|E_G|)-1$ binary AND and OR
gates.  This upper bound is unconditional and is linear on the
bounded-degree family.  It concerns availability of unary choices,
not the full formal polynomial or its integer evaluation.

\paragraph{The classical lower bound being imported.}
Bova--Capelli--Mengel--Slivovsky, \emph{A Strongly Exponential
Separation of DNNFs from CNF Formulas}, arXiv:1411.1995v3
(19 February 2015), define DNNFs and wire size in Section~2,
printed pp.~2--3, and define the graph CNF
\[
                  F_G(x)=\bigwedge_{uv\in E_G}(x_u\vee x_v)
\]
in equation~(2), printed p.~5.  Their Theorem~4, printed p.~10,
states that, for a $(c,d)$ vertex expander,
\[
 \operatorname{DNNF}(F_G)\ge
      2^{g(c,d)\operatorname{size}(F_G)-1},
 \qquad
 g(c,d)=\frac{c\log_2(1+2^{-d})}{6d^3}>0.
\]
Here size counts wires, and the graph CNF has $3|E_G|$ wires.
We import this theorem; its expander lower-bound proof is not
being presented as a new deduction.  The primary manuscript is
\url{https://arxiv.org/pdf/1411.1995v3}.

To relate the hypotheses, the preceding edge-boundary inequality
and maximum degree $d$ give
$|N(S)\setminus S|\ge(h/d)|S|$ for $|S|\le n/2$.
The graph is connected.  A graph of maximum degree at most two
cannot satisfy this fixed positive edge expansion for unbounded
$n$: a connected such graph is a path or cycle, and a consecutive
set of $\lfloor n/2\rfloor$ vertices has at most two boundary edges.
Thus its actual maximum degree $d'$ lies between $3$ and $d$ for
all sufficiently large instances.  Apply the imported theorem
with $c=h/d$ and $d'$.  Since
$g(h/d,d')\ge g(h/d,d)$ and
$3|E_G|\ge3(n-1)\ge n$ for $n\ge2$, we may take
\[
                 \gamma=\frac{h\log_2(1+2^{-d})}{6d^4}.
\]
The finite excluded range has no effect on the asymptotic statement.

\paragraph{The positive-arithmetic-to-DNNF transfer.}
Let a positive arithmetic circuit with $s$ internal binary gates
compute $J_G$ as a formal polynomial.  Delete zero polynomials
and gates not contributing to the output; this is a semantic
simplification used for a lower-bound proof, and no algorithmic
identity test is required.  For each remaining gate $g$, choose
a path to the output and one nonzero monomial from every sibling
at a multiplication on that path.  Their product is a fixed
context monomial $c$.  Positivity ensures that $pc$ is an output
monomial for every monomial $p$ at $g$.  This statement is valid
for a DAG: if another occurrence of a shared gate is in a sibling,
the ordinary distributive expansion permits its monomial to be
chosen independently in that occurrence.

Every output monomial has block multidegree $(1,\ldots,1)$.
Consequently every monomial at $g$ has the same block multidegree,
with entries zero or one.  Denote its block set by $A_g$.
At a productive multiplication gate the two input block sets
are disjoint, and their union is $A_g$.  At an addition gate
the two nonzero input block sets equal $A_g$.
Since coefficients are nonnegative and zero gates have been
removed, the vertex variables appearing anywhere below a gate
are exactly the union of the blocks in its monomials; no
syntactic input can disappear by cancellation.

Replace addition and multiplication by OR and AND, respectively,
each positive constant by true, and the two input rails by
\[
                         z_{v,1}\mapsto x_v,
                 \qquad z_{v,0}\mapsto\neg x_v.
\]
The disjoint block sets just proved show that every AND gate
has variable-disjoint input subcircuits.  The resulting circuit
is therefore a DNNF, with at most $2s$ wires.  On a Boolean
assignment $x$, a monomial of $J_G$ survives this substitution
exactly when its selected vertex set is $\{v:x_v=1\}$.
Thus the DNNF computes the independent-set predicate
$\bigwedge_{uv}(\neg x_u\vee\neg x_v)$.
Interchanging $x_v$ and $\neg x_v$ at all leaves produces a
DNNF for $F_G$ with the same number of wires; decomposability
is unchanged.  The imported theorem now gives
\[
                       2s\ge2^{\gamma n-1},
             \qquad s\ge2^{\gamma n-2},
\]
as asserted.  Arbitrary nonnegative real constants do not
weaken this argument: only zero versus positivity and formal
monomial support were used.

\paragraph{The integral and Boolean upper bounds.}
There are $n$ Boolean variables, one unary factor
\[
                (1-x_v)z_{v,0}+x_vz_{v,1}
\]
for each vertex, and one fixed binary factor $1-x_ux_v$ for
each edge.  Their partition polynomial is exactly $J_G$.
There are at most $(1+d/2)n$ factors of arity at most two,
so NN-C034 applies at the original parameter $n$.  Its formal
integral identity remains valid after this fixed specialization;
no verification variables or growing-scope factors are introduced.
The integral exponent is $o(n)$, whereas the positive lower
exponent is linear, giving an exponential ratio for large $n$.

For equation~\eqref{eq:S019-independent-support}, necessity of
each unary clause follows because every selected assignment
must use an allowed value at that vertex.  If both endpoints
of an edge have their zero values forbidden, both are forced
to one, so its edge clause is necessary too.  Conversely,
suppose all displayed clauses hold.  Choose
$I=\{v:a_{v,0}=0\}$.  The unary clauses ensure that one is
allowed on $I$; zero is allowed everywhere outside $I$ by
definition.  Every edge clause forbids both endpoints from
lying in $I$, so $I$ is independent.  This gives a permissible
monomial and proves the formula.  Its $n+|E_G|$ binary OR
gates and $n+|E_G|-1$ conjunction gates establish the size bound.
This short formula proof concerns the new allowed-weight
interface, not the imported graph-CNF DNNF lower bound.

\paragraph{Geometric and representation statement.}
The integral scheme
\[
 Y_G=\operatorname{Spec}
 \mathbb Z[x_v:v\in V]/
           (x_v^2-x_v, x_ux_v:uv\in E_G)
\]
is the disjoint union of one copy of $\operatorname{Spec}\mathbb Z$
for every independent set.  Put $B=\mathbb Z[z_{v,b}]$.
On the finite locally free $B$-algebra
$B\otimes_{\mathbb Z}\Gamma(Y_G,\mathcal O)$, multiplication
by $\prod_v((1-x_v)z_{v,0}+x_vz_{v,1})$ has trace $J_G(z)$.
This follows directly from its values in the orthogonal
Boolean-idempotent basis.  The trace and arbitrary base-change
interface is EGA~IV4, 18.2.1--18.2.2, printed p.~111, including
equation~18.2.1.2.  The foregoing circuit bound is a uniform
bound for this regular trace function in the compact graph
presentation; it does not assume that the full idempotent
basis is part of the input.

Equivalently, let $T=(\mathbb G_m)^{2n}$ over $\mathbb Z$,
with the specified rail characters.  The actual free
$T$-module
\[
 \mathcal W_G=
 \bigoplus_{I\,\text{independent}}
       \left(\bigotimes_{v\in I}L_{v,1}\right)
          \otimes
       \left(\bigotimes_{v\notin I}L_{v,0}\right)
\]
has character $J_G$.  In the model with these one-dimensional
generators, direct sums, tensor products, nonnegative integer
multiplicities and shared subexpressions, constructing its
class requires $2^{\Omega(n)}$ positive arithmetic gates.
The integral representation Grothendieck ring admits a
$2^{o(n)}$ expression for that same effective class using
virtual differences, by E and NN-C034.  The classical
dictionary is SGA~3, Expos\'e~I, 4.7.2--4.7.3, for torus
gradings, and SGA~6, Expos\'e~IV, (1.5.2) and Definition~2.2,
for the sum, tensor product and Grothendieck-ring operations.
No arbitrary representation functor, entire finite pushforward,
or coefficient-extraction oracle is counted as a unit-cost gate.

For Boolean unary availability data $a$, discard from $Y_G$
the independent-set sections using a forbidden rail.
This is an open-and-closed integral subscheme, whose nonemptiness
is exactly equation~\eqref{eq:S019-independent-support}.
The linear monotone formula thus decides whether an effective
summand is available, while the full character has the proved
exponential positive arithmetic complexity.  Here Boolean
evaluation is the support map on the positive polynomial
semiring; setting a character weight to zero is not asserted
to define a homomorphism of the torus representation ring.
The result distinguishes these explicit arithmetic, Boolean
and representation models.  It gives neither a polynomial
integral circuit nor a lower bound for unrestricted Boolean
circuits, and uses neither A nor M.
